\documentclass[english,twoside, openright, hidelinks, 11 pt, class=report,crop=false]{standalone}
\input{../../preamb}
\input{../bm}


\begin{document}

\grubr{opgderkont}
\abc{
\item Av \eqref{derdefalt} har vi at
\[ f'(a)=\lim\limits_{b\to a} \frac{f(b)-f(a)}{b-a}\]
For at denne grensen skal eksistere, må vi ha at $ \lim\limits_{b\to a} \left(f(b)-f(a)\right)=0 $ (hvis ikke går grensen mot $ \pm \infty $, og da er ikke den deriverte definert), og følgelig er $ \lim\limits_{a\to b} f(b)=f(a) $. Dermed er $ f $ kontinuerlig for alle $ x $.
\item Gitt funksjonen $ f(x)=0 $, da er $ f'(x)=0 $ for alle $ x $. Av resultatet fra a) er dermed $ f(x)=0 $ kontinuerlig. \os

Gitt funksjonen $ g(x)=a $, hvor $ a $ er en konstant. Da er $ f'(x)=0 $, og dermed er $ g(x) $ kontinuerlig.\os

Gitt funksjonene $ h(x)=ax+b $ hvor $ a $ og $ b $ er konstanter. Da er $ h'(x)=g(x) $, og dermed er $ h(x) $ kontinuerlig.

Med samme resonnement kan vi stegvis øke graden av polynomet så høyt vi måtte ønske det, og dermed er alle polynomfunksjoner kontinuerlige.
}

\grubr{opgder2dx} \\

\abc{ \item Vi har at
\alg{
\lim\limits_{h\to0}\frac{f(x+h)-f(x-h)}{h}&=\lim\limits_{h\to0}\frac{f(x+h)-f(x)+f(x)-f(x-h)}{h} \\
&= \lim\limits_{h\to0}\frac{f(x+h)-f(x)}{h}+\lim\limits_{h\to0}\frac{f(x)-f(x-h)}{h} \\
&= \lim\limits_{h\to0}\left[f'(x)+f'(x-h)\right] \\
&= 2f'(x)
}

\newpage
\grubr{opgderarchi}
\fig{opgderarchi_lf}
\abc{
\item Vi har at
\[ A=(p, ap^2+ bp+c)\qquad,\qquad B = (q, aq^2+bq+c) \]
}
Vi lar $g(x)$ være funksjonen som beskriver linja gjennom $AB$. Da er
\alg{
 g'(x)&= \frac{f(q)-f(p)}{q-p} \br
 &=\frac{aq^2+bq+c-(ap^2+ bp+c)}{q-p} \br
 &=\frac{a(q^2-p^2)+b(q-p)}{q-p} \br
 &=\frac{a(q+p)(q-p)+b(q-p)}{q-p} \br
 &=a(q+p)+b
}
Vi setter $C=(t, f(t))$. Da har vi at
\alg{
f'(t)&=a(q+p)+b \\
2a t+ b &= a(q+p)+b \\
t &= \frac{1}{2}(q+p)
}
\item Vi lar $D$ være midtpunktet til $AB$. Da er $D=\left(t, \frac{1}{2}(f(q)+f(p))\right)$ (se \refops{vekopgmidt}). Videre er
\alg{
DC&=\frac{1}{2}\left(f(q)+f(p)\right)-f(t) \br
&= \frac{1}{2}\left(a(p^2+q^2)+b(p+q)+2c\right)-a\left(\frac{1}{2}(p+q)\right)^2+b \left(\frac{1}{2}(p+q)\right)+c \br
&= \frac{a}{2}(p^2+q^2)-\frac{a}{4}(p+q)^2 \br
&=\frac{a}{4}(2p^2+2q^2-p^2-q^2-2pq) \\
&=\frac{a}{4}(q-p)^2
}
}
Vi plasserer $E$ og $F$ slik at $\square AEFB$ er et parallellogram. Med $AE$ som grunnlinje, har $\square AEFB$ høyde $q-p$. Da $AE=DC$, er
\nn{
A_{\triangle ABC}=\frac{1}{2}A_{\square AEFB}=\frac{1}{2}\cdot\frac{a}{4}(q-p)^2(q-p)=\frac{a}{8}(q-p)^3
}


\opr{opgderdernotdefnd}
Vi har at
\[ \lim\limits_{h\to0^-} \frac{f(h)-f(0)}{h}=\lim\limits_{h\to0^-} \frac{g(h)-g(0)}{h}=g'(0)=1 \]
\begin{align*}
	\lim\limits_{h\to0^+} \frac{f(h)-f(0)}{h}&=\lim\limits_{h\to0^+} \frac{k(h)-g(0)}{h} \\
	&= \lim\limits_{h\to 0^+} \frac{h^2+h-1}{h} \br
	&=-\infty
\end{align*}
Dermed eksisterer ikke $f'(0)$.
\newpage
%\grubr{directrix}
\textbf{Gruble 27}\\
\fig{directrix_lf}
\abc{
\item Av $y$-verdien til $A$ har vi at
\begin{align*}
&&	p&=f(2p) \\
&&	p&= 4p^2 && (p\neq 0) \\
&& 1 &= 4p \\
&& p &= \frac{1}{4}
\end{align*}
\item Av \pyt på $\triangle AKC$ har vi at
\alg{
&&FC^2 &= FK^2+CK^2 \\
&&&= c^2+(f(c)-p)^2 \\
&&&= c^2 + (c^2-p)^2 \\
&&&=c^2 +c^4 -2c^2p+p^2 \\
&&&= c^2+c^4-\frac{1}{2}c^2+\frac{1}{16} \\
&&&=c^4+\frac{1}{2}c^2+\frac{1}{16} \\
&&&= \left(c^2+\frac{1}{4}\right)^2 && (FC>0) \\
&&FC&= c^2+\frac{1}{4} 
}
\item 
Vi har at
\[ CD=CJ+JD=f(c)+p=c^2+\frac{1}{4}=FC \]
\item Vi har at $f'(c)=2c$. Tangenten til $f$ i $C$ kan derfor beskrives ved funksjonen $g(x)=2cx+b$. I punktet $C$ har vi at
\begin{align*}
g(c) &= f(c) \\
2c^2+b &= c^2 \\
b &= -c^2
\end{align*}
I punket $E$ har vi at
\alg{
&&g(x) = 0 && \\
&&2cx-c^2 = 0 && (c\neq0) \\
&& 2x-c &= 0 \\
&& x &= \frac{c}{2}
}
$\triangle FOE\cong\triangle DJE$ fordi begge er rettvinklede, $OF=JD=p$ og $OE=EJ=\dfrac{c}{2}$. Følgelig er $E$ midptunktet på $FD$.
\item Da $\triangle FCD$ er likebeint og $E$ er midtpunktet på $FD$, er $\angle ACE=\angle ECD$. $\angle HCG=\angle ECD$ fordi de er toppunkt. Dermed er $\angle FCE=\angle HCG$.
}

\newpage
\grubr{opgderpointb}
\fig{opgderpointb_lf}

\abc{
\item Da $\vec{v}$ danner vinkelen $15^\circ$ med $x$-aksen, er $[1, \tan 15^\circ]$ en retningsvektor for $\vec{v}$. Dermed er
\[ \vec{v}(t)=O+t[1, \tan 15^\circ]=t[1, 2-\sqrt{3}] \]
\item Vi har at
\begin{align*}
	\vv{BO}&=-t[1, 2-\sqrt{3}] \vn
	\vv{BC}&=[4\sqrt{3}-5-t, -1-t(2-\sqrt{3})]
\end{align*}
Følgelig er
\begin{align*}
	\vv{BO}\cdot\vv{BC}&=-(4\sqrt{3}-5)t+t^2+(2-\sqrt{3})t+(2-\sqrt{3})^2t^2 \\
	&=(8-4\sqrt{3})t^2+(7-5\sqrt{3})t
\end{align*}
Av cosinusstetningen har vi at
\begin{align*}
	|BO||BC|&=-\frac{|OC|^2-|BO|^2-|BC|^2}{2\cos \angle OBC} \br
	&=-\frac{24(7\sqrt{3}-13)}{2\cos \angle OBC} \br
	&=\frac{12(13-7\sqrt{3})}{\cos \angle OBC}
\end{align*}
Nå har vi at
\begin{align*}
	\frac{\vv{BO}\cdot\vv{BC}}{|BO||BC|}&=\cos \angle OBC \br
	\left((8-4\sqrt{3})t^2+(7-5\sqrt{3})t\right) \frac{\cos \angle OBC}{12(13-7\sqrt{3})} &= \cos \angle OBC \\
	(8-4\sqrt{3})t^2+(7-5\sqrt{3})t &= 12(13-7\sqrt{3}) \\
	(8-4\sqrt{3})t^2+(7-5\sqrt{3})t - 12(13-7\sqrt{3})&=0
\end{align*}
Dermed er
\[ f(t)=(8-4\sqrt{3})t^2+(7-5\sqrt{3})t - 12(13-7\sqrt{3}) \]
\item Ekstremalverdien til $f$ er gitt ved
\begin{align*}
	f'(t)&=0 \\
	2(8-4\sqrt{3})t+(7-5\sqrt{3}) &=0 \\
	8(2-\sqrt{3})t&=5\sqrt{3}-7 \\
	t &= \frac{5\sqrt{3}-7}{8(2-\sqrt{3})} \br
	&= \frac{(5\sqrt{3}-7)}{8(2-\sqrt{3})}\cdot \frac{2+\sqrt{3}}{2+\sqrt{3}} \br
	&=\frac{10\sqrt{3}-14+15-7\sqrt{3}}{8} \br
	&=\frac{3\sqrt{3}+1}{8}
\end{align*}
\item Vi vet at ekstremalpunktet ligger midt mellom nullpunktene til en andregradsfunksjon. Avstanden mellom ekstremalpunktet og det gitte nullpunktet er
\[\frac{3\sqrt{3}+1}{8}-\frac{3\sqrt{3}-15}{4}=\frac{31-3\sqrt{3}}{8}\]
Det andre nullpunktete er dermed gitt ved
\[ \frac{3\sqrt{3}+1}{8}+\frac{31-3\sqrt{3}}{8}=4 \]
I punktet $B$ er dermed $t=4$, og da er
\[ B=\vec{v}(4)=(4, 4(2-\sqrt{3})) \]
}
\newpage

\end{document}